\documentclass{phys400finalreport}
\usepackage[dvipsnames]{xcolor}
\usepackage{bm}

\proposaltitle{Quantifying Information Extracted from
Gravitational Wave Observations: An Empirical Test of the
Flanagan--Hughes Framework Using GWTC Posterior Samples}
\studentname{Ömer Faruk Soydan}
\studentid{2555472}
\advisor{Assist. Prof. Dr. Doğa Veske}
\studyfield{Gravitational Wave Astrophysics}
\studytype{Computational Data Analysis, Theoretical Testing}
\submissiondate{\today}

\begin{document}
\makeproposalcover
\newpage

%==================================================================
% 1. INTRODUCTION
%==================================================================
\introduction

The direct detection of gravitational waves (GWs) has opened a fundamentally
new observational window into the universe \citep{Abbott2016}. The
LIGO--Virgo--KAGRA (LVK) collaboration has completed four observing runs since
the first observation of GW150914. The most recent data release, GWTC-5.0,
brings the total number of compact binary merger candidates in the cumulative
catalog to 390
\citep{theligoscientificcollaboration2026gwtc50observationssecondfourth}. Each
detected event carries physical information about its
source, including the component masses, spins, and sky location. This
information is extracted through Bayesian parameter estimation (PE), which is
implemented in dedicated samplers such as those in LALInference and its
successors \citep{Veitch2015}. The resulting posterior samples are made
publicly available through the Gravitational Wave Open Science Center
\citep{ligo_scientific_collaboration_and_virgo_2023_8177023} and are read
using the \texttt{pesummary} library \citep{Hoy2021}. With next-generation
detectors on the horizon, such as the Einstein Telescope and Cosmic Explorer,
understanding the information-theoretic limits of these observations has become
increasingly important.

The amount of information extracted by a gravitational wave observation can be
quantified rigorously using classical information theory. \citet{Shannon1948}
showed that the reduction in uncertainty produced by a measurement can be
expressed in units of bits. In a Bayesian context, this quantity is the
\citet{1320776d-9e76-337e-a755-73010b6e4b64} (KL) divergence from the posterior
to the prior distribution. The amount of information that can be extracted
about a gravitational wave source's parameters is given by the following exact
expression:
\begin{equation}
I_\mathrm{source} = \int d\theta\, p(\theta | s) \log_2\left(\frac{p(\theta | s)}{p(\theta | p(s))}\right)
\label{eq:exact_info}
\end{equation}
where $p(\theta | s)$ is the posterior distribution given the observed signal
and $p(\theta | p(s))$ is the prior probability distribution given the
magnitude of the measured signal $p(s)$. Computing this integral directly from
real GWTC posterior samples is the central goal of this project.

\citet{Flanagan1998} introduced an analytical estimate of this quantity by
assuming that detector noise is Gaussian and stationary, and that the waveform
$h(\theta)$ depends linearly on the source parameters $\theta$, where there is
little prior information. Under these conditions, Equation~\ref{eq:exact_info}
reduces to:
\begin{equation}
I_\mathrm{source} \approx \frac{N_\mathrm{param}}{2} \log_2\left(1 + \frac{\rho^2}{N_\mathrm{param}}\right)
\label{eq:approx_info}
\end{equation}
where $\rho$ is the matched-filter SNR and $N_\mathrm{param} = 15$ for a binary
black hole system, which covers the component masses, spin magnitudes and
orientations, and extrinsic parameters. For a GW150914-like system this
predicts approximately $41.5$ bits. The approximation performs well in the
high-SNR regime, but it becomes less reliable as the SNR decreases, because
non-Gaussian features in the posterior become more prominent. Despite its
elegance, this framework has not been systematically tested against the actual
posterior distributions produced by modern PE pipelines; the original authors
themselves describe it as ``a rather crude approximation.''

This project aims to address this gap. The central research question is: of all
the gravitational wave events detected to date, how much information about
source parameters has been obtained, and how closely does this agree with the
prediction of Flanagan and Hughes? We hypothesise that the measured information
will be close to, but not exactly equal to, the analytical prediction. This
departure is expected to be most noticeable for parameters that are strongly
correlated or poorly constrained, and at low signal-to-noise ratio (SNR), where
non-Gaussian features are most pronounced.

To answer this question, we developed a Python-based pipeline that evaluates
the KL divergence directly from GWTC posterior samples over the full
15-dimensional parameter space of a binary black hole merger. Because direct
density estimation in 15 dimensions is obstructed by the curse of
dimensionality, the analysis combines several complementary estimators, rather
than relying on a single method, as described in Section~\ref{sec:methods}. The
pipeline is first validated on the benchmark event GW150914 and then applied to
a sample of 35 binary black hole mergers, so that the Flanagan--Hughes
prediction can be tested both for a single event and across the population. The
remainder of this report is organised as follows. Section~\ref{sec:methods}
describes the data and the estimation methods, Section~\ref{sec:results}
presents and discusses the single-event and population results, and
Section~\ref{sec:conclusion} summarises the main conclusions.

%==================================================================
% 2. METHODS
%==================================================================
\methods\label{sec:methods}

\subsection*{Data}

The analysis uses the publicly available parameter-estimation (PE) samples
from the IGWN / Gravitational Wave Open Science Center (GWOSC) GWTC data
release. The sample comprises 35 binary black hole (BBH) events: GW150914 from
GWTC-2.1 together with 34 events from GWTC-3, read directly from the public
HDF5 products. For each event the posterior is taken from the
\texttt{C01:IMRPhenomXPHM} waveform model and contains of order $10^5$ samples,
ranging from a few thousand to roughly $2.5 \times 10^5$, while the prior is
the smaller distribution (about 5000 samples) stored alongside it. The samples
are read with the \texttt{pesummary} library \citep{Hoy2021}, with a direct
\texttt{h5py} fallback, and repeated posterior rows are removed before
analysis. The 15 source parameters are the masses, spin magnitudes and tilts,
azimuthal spin angles, luminosity distance, inclination, polarization, sky
position, coalescence time, and phase. The geocentric time $t_c$ is treated on
a relative scale, because its absolute value of order $10^9$~s would otherwise
dominate the distance metrics used below.

\subsection*{Information Measure and the Marginal Lower Bound}

The information gained by an observation is the KL divergence of
Equation~\ref{eq:exact_info}, expressed in bits. To validate the pipeline, the
one-dimensional (1D) KL divergence of the primary mass $m_1$ of GW150914 was
computed by two independent routes: kernel density estimation (KDE) and a
binned entropy calculation. The two agreed to within about $0.01$~bit
($2.638$ versus $2.626$~bits), confirming that the density estimation and
integration steps were correct.

Summing the 1D KL divergences of all 15 parameters yields a marginal total.
This sum is not the full information, because the parameters are correlated.
The joint divergence decomposes as
\begin{equation}
I_\mathrm{joint} = \sum_i I_\mathrm{marg}(i) + \mathrm{TC},
\label{eq:tc}
\end{equation}
where $\mathrm{TC}$ is the total correlation, which is always non-negative. The
marginal sum therefore provides a guaranteed lower bound on the joint
information; for GW150914 it is approximately $37$~bits. Recovering the
remaining contribution requires estimating the joint 15-dimensional density.

\subsection*{The Curse of Dimensionality}

Direct density estimation in 15 dimensions is severely constrained, as
summarised in Table~\ref{tab:methods}. A direct histogram
(\texttt{np.histogramdd} \citep{numpy_paper}) with only five bins per axis
would require $5^{15} \approx 3 \times 10^{10}$ cells, demanding on the order of
$240$~GB of memory and leaving almost every cell empty for $\sim\!10^5$
samples; it is therefore infeasible \citep{altman2018curse}. A full
15-dimensional $k$-nearest-neighbour ($k$-NN) estimator does not converge: its
value drifts from $25.9$ to $104.9$~bits as the sample size grows, because
Euclidean distances become poorly calibrated in high dimensions. A multivariate
Gaussian model \citep{munoz2026kl} gives a fast closed-form result
($\approx\!41.2$~bits for GW150914, close to the analytical prediction), but it
assumes a Gaussian posterior and produces an ill-conditioned covariance matrix
(condition number $\sim\!10^9$). The common cause of all these failures is
dimensionality.

\begin{table}[htbp]
\centering
\small
\caption{Comparison of candidate estimators for the 15-dimensional KL
divergence, evaluated on GW150914. The first four entries illustrate the
difficulty of direct high-dimensional estimation; Group-KLD is the method
adopted in this work.}
\label{tab:methods}
\begin{tabular}{lllc}
\toprule
Estimator & Approach & GW150914 & Status \\
\midrule
Marginal 1D sum (KDE)        & Sum of independent 1D KL divergences        & $\approx 37$ bit       & Lower bound only \\
\texttt{np.histogramdd}      & Direct $n$-dimensional grid count           & ---                    & Infeasible ($\sim$240 GB) \\
Full 15D $k$-NN              & Neighbour-distance density ratio            & $25.9 \rightarrow 104.9$ bit & Does not converge \\
Multivariate Gaussian        & Closed-form covariance log-determinant      & $\approx 41.2$ bit     & High-SNR baseline; ill-conditioned \\
\textbf{Group-KLD (adopted)} & Five $\leq$5D blocks, 3 estimators, TC corr. & $\approx 39$ bit       & Method-independent \\
\bottomrule
\end{tabular}
\end{table}

\subsection*{Cross-Checking Against a Published Benchmark}

Because no single estimator can be trusted in many dimensions, the choice of
method was guided by the benchmark of \citet{alvarezchaves2024}, who tested
KDE, binning, and nearest-neighbour estimators against analytic ground truth in
up to ten dimensions. In their results the nearest-neighbour estimator
converges for the KL divergence, KDE drifts slowly, and binning fails outright.
Two design rules follow directly: cross-check several estimators rather than
rely on one, and keep every density estimate to at most five dimensions, where
all of these methods remain reliable.

\subsection*{Group-KLD with a Hybrid Prior}

The method adopted in this work, here called Group-KLD, applies both rules. The
15-parameter vector is split into five physically motivated blocks --- the two
masses; the spin magnitudes and tilts; the azimuthal spin angles; the distance
and orientation; and the sky position, time, and phase --- each of at most five
dimensions. The blocks are chosen so that correlations within a block are
strong while correlations between blocks are weak. Within each block the KL
divergence is computed with three independent estimators, KDE-Scott
\citep{scott1992}, KDE-Silverman \citep{silverman1986}, and a $k=1$
nearest-neighbour estimator \citep{Perez-Cruz2008, Wang2009}, and their mutual
agreement is used as a consistency check.

Two corrections complete the estimate. First, a hybrid prior addresses cases
where the posterior extends beyond the support of the stored prior. When a
posterior overflows the stored prior by more than $0.5\%$, which occurs most
strongly for the luminosity distance and the masses, the stored prior is
replaced by an analytic prior for that parameter; otherwise the original prior
is bootstrap-resampled. This removes a spurious contribution that would
otherwise arise from the limited support of the stored prior samples. Second, a
total-correlation correction restores the between-block dependence that the
block decomposition would otherwise omit:
\begin{equation}
I_\mathrm{joint} \approx \sum_{g} I_\mathrm{group}(g) + \left(\mathrm{TC}_\mathrm{post} - \mathrm{TC}_\mathrm{prior}\right),
\label{eq:group_kld}
\end{equation}
where the sum runs over the five blocks. The resulting value is the joint
information estimate used throughout this report. For GW150914 the three
estimators agree at $\approx\!39$~bits, and the identical pipeline is applied,
without modification, to all 35 events.

%==================================================================
% 3. RESULTS AND DISCUSSION
%==================================================================
\resultsanddiscussion\label{sec:results}

\subsection*{Single-Event Information: GW150914}

Applied to the benchmark event GW150914, the Group-KLD method returns a joint
information of approximately $39$~bits. The three independent estimators
(KDE-Scott, KDE-Silverman, and the nearest-neighbour estimator) agree closely,
both in the total and in the per-block contributions, as shown in
Figure~\ref{fig:groupkld}. The agreement confirms that the result is
method-independent and not an artefact of any single density estimator. The
information is dominated by the blocks containing the masses, the distance and
orientation, and the sky position and time, while the spin-angle blocks
contribute little, consistent with these parameters being only weakly
constrained for this event.

\begin{figure}[h]
\centering
\includegraphics[width=0.75\textwidth]{grup_kld_hibrit_GW150914.png}
\caption{Per-block and total Group-KLD information for GW150914, computed with
three independent estimators (KDE-Scott, KDE-Silverman, and a $k=1$
nearest-neighbour estimator) using the hybrid prior. The three estimators
overlap in every block, giving a method-independent total of $\approx\!39$~bits.}
\label{fig:groupkld}
\end{figure}

This value can be placed against the other estimators of Table~\ref{tab:methods}.
The measured information orders as expected: the marginal sum
($\approx\!37$~bits) is a strict lower bound, the correlation-aware Group-KLD
estimate ($\approx\!39$~bits) lies above it, and both fall below the
multivariate Gaussian baseline ($\approx\!41.2$~bits) and the Flanagan--Hughes
prediction ($\approx\!41.5$~bits). The ordering
$I_\mathrm{marg} \le I_\mathrm{joint} \lesssim I_\mathrm{Gauss} \approx I_\mathrm{FH}$
is exactly what the decomposition of Equation~\ref{eq:tc} requires. The
Flanagan--Hughes value is therefore an accurate but slightly high baseline: the
direct measurement sits about two bits lower, a gap that reflects the genuine
non-Gaussianity of the posterior, which the Gaussian and analytical estimates
do not capture.

\subsection*{Population Scaling Across 35 Events}

The same pipeline, applied unchanged to all 35 events, yields information gains
that span $6.7$~bits for the quietest event (GW200322) to $44.6$~bits for the
loudest (GW191216). This range is not random: the information grows
systematically with the loudness of the signal. Figure~\ref{fig:scaling} shows
the joint information against the network matched-filter SNR, together with
three candidate fits. A logarithmic law, $I = a\ln(\mathrm{SNR}) + b$, gives the
best single-variable description (Pearson $r = 0.72$), and outperforms both a
linear and a power-law fit, as compared in Table~\ref{tab:fits}.

\begin{figure}[h]
\centering
\includegraphics[width=0.8\textwidth]{result_test_plot.png}
\caption{Joint information gain against the median network matched-filter SNR
for the event sample. The logarithmic fit $I = a\ln(\mathrm{SNR}) + b$
(solid line) describes the trend better than the linear and power-law
alternatives (dashed lines). Information increases with signal loudness, as
expected when it scales with the square of the SNR.}
\label{fig:scaling}
\end{figure}

\begin{table}[htbp]
\centering
\caption{Fits of the joint information $I$ (in bits) across the event sample.
The logarithmic dependence on SNR is the best single-variable model; adding the
total mass $M$ and the number of detectors $n_\mathrm{det}$ raises the explained
variance to $0.88$.}
\label{tab:fits}
\begin{tabular}{lc}
\hline
Model & $R^2$ \\
\hline
$I = m\,\mathrm{SNR} + c$ \quad (linear)        & 0.52 \\
$I = A\,\mathrm{SNR}^{p}$ \quad (power law)      & 0.54 \\
$I = a\ln(\mathrm{SNR}) + b$ \quad (logarithmic) & 0.58 \\
$I = a\ln(\mathrm{SNR}) + c\ln(M) + d\,n_\mathrm{det}$ & 0.88 \\
\hline
\end{tabular}
\end{table}

The logarithmic dependence is not merely an empirical fit; it is precisely the
behaviour predicted by Flanagan and Hughes. In the high-SNR limit,
Equation~\ref{eq:approx_info} reduces to
$I \approx N_\mathrm{param}\log_2(\mathrm{SNR})$, so the information should grow
as the logarithm of the SNR. The data confirm this functional form across the
whole sample, even though the absolute Gaussian assumption is imperfect for
individual events. In other words, the Flanagan--Hughes framework captures
\emph{how} information scales with loudness, while the direct measurement is
needed to pin down its absolute value event by event.

Signal loudness alone, however, explains only part of the variance
($R^2 = 0.58$). Two further physical factors close most of the remaining gap.
Adding the total mass and the number of participating detectors gives the
three-term law
\begin{equation}
I = a\ln(\mathrm{SNR}) + c\ln(M) + d\,n_\mathrm{det},
\label{eq:fulllaw}
\end{equation}
which explains $88\%$ of the variance. Each term has a clear physical meaning.
The coefficient of $\ln(\mathrm{SNR})$ is positive: louder signals carry more
information. The coefficient of $\ln(M)$ is negative: heavier binaries reveal
less, because lighter binaries remain in the detectors' sensitive band for many
more cycles and therefore imprint more measurable structure. The coefficient of
$n_\mathrm{det}$ is positive: additional detectors tighten the sky-position and
distance measurements, adding information that a single detector cannot supply.

\subsection*{Limitations}

Three limitations qualify these results. First, although the $\leq$5-dimensional
blocking keeps each density estimate in a regime where the methods agree, the
between-block correlation is only approximately recovered through the
total-correlation correction of Equation~\ref{eq:group_kld}; a fully direct
15-dimensional joint estimate remains computationally out of reach. Second, the
stored prior samples do not always cover the full extent of the posterior,
most notably in luminosity distance and the masses. The hybrid prior corrects
this where it occurs, but it remains the largest source of systematic
uncertainty, and an uncorrected Gaussian or nearest-neighbour estimate inflates
the information in exactly these directions. Third, the population fit uses the
events for which a usable network SNR could be recovered, and the scatter about
the single-variable trend is genuine, reflecting differences in mass, detector
network, and non-Gaussian structure between events rather than measurement
noise alone.

%==================================================================
% 4. CONCLUSION
%==================================================================
\conclusion\label{sec:conclusion}

This project set out to measure, in bits, how much information about source
parameters a gravitational wave detection actually provides, and to test how
well the classic analytical estimate of Flanagan and Hughes survives when
confronted with real posterior distributions. The information was defined as
the Kullback--Leibler divergence from the posterior to the prior and computed
directly from public GWTC parameter-estimation samples. Because no single
estimator is reliable in the full 15-dimensional parameter space, a grouped
method (Group-KLD) was developed, in which the parameters are split into five
physically motivated blocks of at most five dimensions, a hybrid prior removes a
stored-prior support artefact, and three independent estimators are
cross-checked for agreement.

The research question can now be answered directly. For the benchmark event
GW150914 the measured information is approximately $39$~bits, about two bits
below the Flanagan--Hughes prediction of $41.5$~bits. Across the full sample of
35 binary black hole events the information ranges from $6.7$ to $44.6$~bits.
These two findings together show that the Flanagan--Hughes framework is an
accurate but slightly optimistic baseline: it reproduces the correct order of
magnitude and, more importantly, the correct functional dependence, while
overstating the absolute value by a small margin that is traceable to the
non-Gaussianity of real posteriors.

The central result of the population analysis is a clean scaling law. The
information grows as the logarithm of the signal-to-noise ratio, exactly as the
high-SNR limit of the Flanagan--Hughes formula predicts, and a three-term law
combining loudness, total mass, and the number of participating detectors
explains $88\%$ of the variance across the catalogue. Louder, lighter binaries
observed by more detectors are the most informative: louder signals carry more
information, lighter binaries remain in band for more cycles, and additional
detectors sharpen the sky-position and distance measurements. This law
summarises, in a single compact expression, how efficiently the present
LIGO--Virgo--KAGRA network converts a detection into knowledge about its source.

The significance of these results is twofold. First, they establish an
information-theoretic benchmark for the current generation of detectors, against
which the gains promised by next-generation observatories such as the Einstein
Telescope and Cosmic Explorer can be quantified. Second, they demonstrate that
a long-standing analytical estimate, while functionally correct, must be
corrected for non-Gaussianity and finite prior support if it is to give
accurate per-event numbers, and they provide a method that does so at the scale
of a full catalogue.

Several directions would extend this work. The pipeline is now fast enough to
be applied to the complete GWTC-5.0 release, enlarging the sample well beyond
the 35 events analysed here and tightening the fitted scaling law. The same
framework could be extended to binary neutron star and neutron star--black hole
systems, whose posteriors differ in character, and to a population-level
information measure that treats the catalogue as a whole rather than event by
event. Finally, neural estimators of mutual information offer a possible route
to a fully joint 15-dimensional estimate that does not rely on the block
decomposition, should the curse of dimensionality be overcome.

\referencessection
\end{document}
